\scanpage{265}
\mspage{12}
\margin{n° 227}
normale) la fonction $f_i z$ est définie à chaque point de $V_i$, et induit donc une fonction $f'_i z'$ sur $W$ ou $W^*$, définie à chaque point de $W$ ayant un représentant sur $V_i$. Si $V'_i$ est le représentant affine de $W$ dérivé de $V_i$, on voit par ce qui précède que les pôles de $z'$ ayant un représentant sur $V'_i$ sont bornés par les zéros de $f'_i$. La variété $W$ étant couverte par un nombre fini de ces variétés affines, on voit que les pôles de $z'$ sont uniformément bornés, cqfd.

Revenons à la démonstration du théorème. Sachant que $L(0)=k$, il suffit de montrer que si le théorème est vrai pour un diviseur positif $D$, alors pour tout diviseur premier $W$, $L(D+W)$ est de dimension finie.

Supposons que $W$ apparaisse avec la multiplicité $n$ dans $D$, i.e. $D=nW+\ldots$. S'il n'y a pas de fonction $f\in L(D+W)$ ayant un pôle d'ordre $n+1$ sur $W$, alors $L(D+W)=L(D)$, et nous avons fini. Autrement, supposons qu'il existe une fonction $f$ dans $L(D+W)$ ayant un pôle d'ordre $n+1$ sur $W$. Considérons l'espace vectoriel $f^{-1}L(D+W)$, formé de tous les produits $f^{-1}g$, avec $g\in L(D+W)$. Il est évidemment égal à $L(D+W+(f))$. Aucune des fonctions de cet espace n'a de pôle sur $W$. Par conséquent chacune de ces fonctions induit une fonction sur $W$, et on a un homomorphisme $h\longrightarrow h'$ où $h\in L(D+W+(f))$, et $h'$ est la fonction induite sur $W$. Le lemme, et l'hypothèse d'induction, montrent que l'image de notre homomorphisme est de dimension finie.

Si une fonction $f^{-1}g$ induit la constante $0$ sur $W$, alors $f^{-1}g$ a un zéro sur $W$, et $g$ a donc un pôle d'ordre au plus $n$ sur $W$. Mais $g$ étant dans $L(D+W)$, on voit que $g$ est même dans $L(D)$. Le noyau de notre homomorphisme est donc précisément l'espace $L(D)$, qui est

\mspage{264}
